\documentclass[11pt]{article}
\usepackage[margin=1in]{geometry}
\usepackage{amsmath,amssymb,amsthm}
\usepackage[hidelinks]{hyperref}

\newtheorem*{theorem}{Theorem}
\newtheorem*{lemma}{Lemma}
\newcommand{\C}{\mathbb C}
\newcommand{\Z}{\mathbb Z}

\begin{document}

\begin{center}
{\Large Literature-implied answer for arXiv:2312.05617}\\[0.5em]
{\large The infinite-dihedral square has the SOS property}
\end{center}

\paragraph{Status.}
\textbf{Literature-implied answer, full scope for the source question.}
Mehta--Slofstra--Zhao ask in the introduction after Corollary 1.2 whether
\(\C\Z_2^{*2}\times \Z_2^{*2}\) has the SOS property; at the end of Example 6.8
they equivalently say that, for \(G=\Z_2^{*2}\), they do not know whether
\(\C G\times G\) is archimedean closed. The answer is yes.

\paragraph{Supporting theorems.}
Dritschel's two-variable matrix Fejer--Riesz theorem says that every positive
finite-dimensional operator-valued trigonometric polynomial in two variables is
a finite sum of hermitian squares of analytic matrix polynomials
\cite[Introduction theorem]{Dritschel2024}. Equivalently, for every \(n\),
the type-I Positivstellensatz holds for
\[
M_n\bigl(\C[z^{\pm1},w^{\pm1}]\bigr)
\]
with positivity understood pointwise on \(\mathbb T^2\).

Savchuk--Schmudgen prove that if the type-I Positivstellensatz holds for
\(M_n(A)\), then it also holds for the finite crossed product
\(A\times_\alpha F\), where \(F\) is a finite group of \(*\)-automorphisms of
\(A\) and \(n=|F|\). Their proof embeds the crossed product into \(M_n(A)\)
by the regular covariant representation and applies a strong conditional
expectation back to the crossed product \cite[Section 5, Theorem 3]{SS2010}.

\begin{theorem}
Let \(G=\Z_2^{*2}\) be the infinite dihedral group. Then every positive element
of the group algebra \(\C[G\times G]\) is a finite sum of hermitian squares in
\(\C[G\times G]\). Hence \(\C\Z_2^{*2}\times \Z_2^{*2}\) has the SOS property,
or equivalently is archimedean closed.
\end{theorem}

\begin{proof}
Write \(G=\langle r,s\mid s^2=1,\;srs=r^{-1}\rangle\cong \Z\rtimes C_2\).
Then
\[
G\times G\cong \Z^2\rtimes F,\qquad F=C_2\times C_2,
\]
where the two generators of \(F\) invert the two coordinates of \(\Z^2\).
Consequently
\[
\C[G\times G]\cong
\C[z^{\pm1},w^{\pm1}]\rtimes_\alpha F.
\]

Let \(A=\C[z^{\pm1},w^{\pm1}]\). By Dritschel's theorem, every pointwise
positive element of \(M_4(A)\) is a sum of hermitian squares in \(M_4(A)\).
Apply Savchuk--Schmudgen's finite crossed-product transfer theorem with
\(|F|=4\). It follows that every element of \(A\rtimes_\alpha F\) that is
positive in the regular covariant representation is a sum of hermitian squares
inside \(A\rtimes_\alpha F\).

It remains only to match this positivity with the representation-positive
notion used by Mehta--Slofstra--Zhao. The action of the finite group \(F\) on
\(\mathbb T^2\) gives the faithful regular representation of the full crossed
product \(C(\mathbb T^2)\rtimes F\) into \(M_4(C(\mathbb T^2))\); equivalently,
the canonical conditional expectation onto \(C(\mathbb T^2)\) is faithful.
Since \(F\) is finite, these regular covariant fibres contain the finite
isotropy representations at fixed points, and positivity in the faithful
regular representation is exactly positivity in the full crossed-product
\(C^*\)-algebra. Restricting to the algebraic crossed product gives the
required representation-positive notion for \(\C[G\times G]\).
\end{proof}

\paragraph{Search evidence.}
Targeted searches for the exact phrases ``\(C Z_2^{*2}\) SOS property'',
``archimedean closed \(Z_2^{*2}\)'', and ``infinite dihedral sums of squares
group algebra'' did not find a paper explicitly stating that the Mehta--
Slofstra--Zhao endpoint question is resolved. The answer above is therefore
recorded as an agent-identified implication of known theorems rather than as a
paper whose authors explicitly advertised the resolution.

\paragraph{Scope.}
This note answers the source's infinite-dihedral-square SOS/archimedean-closed
question. It does not address the \(\Pi^0_2\)-completeness conjecture for the
larger free algebras, the tracial positivity conjecture, or smallest generator
bounds for tracial positivity.

\begin{thebibliography}{9}

\bibitem{MSZ2026}
Arthur Mehta, William Slofstra, and Yuming Zhao.
\newblock \emph{Positivity is undecidable in tensor products of free algebras}.
\newblock arXiv:2312.05617v2; Journal of the European Mathematical Society,
online first, 2026.

\bibitem{SS2010}
Yurii Savchuk and Konrad Schmudgen.
\newblock \emph{Positivstellensatze for Algebras of Matrices}.
\newblock arXiv:1004.1529v2; Linear Algebra and its Applications
436 (2012), 758--788.

\bibitem{Dritschel2024}
Michael A. Dritschel.
\newblock \emph{Factoring non-negative operator valued trigonometric
polynomials in two variables}.
\newblock Mathematische Annalen 391 (2025), 515--537; published online 2024.

\end{thebibliography}

\end{document}
